%==============================================================================
%  A Bias-Aware Virtual Patient Simulator
%  B.Tech Summer Internship 2026, DAU (Dhirubhai Ambani University)
%
%  COMPILE: upload to https://overleaf.com, set compiler to pdfLaTeX, Recompile.
%
%  All results are real:
%    Table II  : from validate_detectors.py
%    Table III : counts from cases.py
%    Table IV,V: from analyze_sessions.py on our 16 collected sessions
%==============================================================================
\documentclass[conference]{IEEEtran}
\IEEEoverridecommandlockouts

\usepackage{cite}
\usepackage{amsmath,amssymb,amsfonts}
\usepackage{graphicx}
\usepackage{booktabs}
\usepackage{array}
\usepackage{url}
\usepackage[hidelinks]{hyperref}

\begin{document}

\title{A Bias-Aware Virtual Patient Simulator with\\
Reasoning-Process Feedback for Medical Education}

\author{
\IEEEauthorblockN{Vraj Patel}
\IEEEauthorblockA{\textit{Roll No. 202301408} \\
\textit{DAU (Dhirubhai Ambani University)}\\
Gandhinagar, India}
\and
\IEEEauthorblockN{Yogesh Bagotia}
\IEEEauthorblockA{\textit{Roll No. 202301114} \\
\textit{DAU (Dhirubhai Ambani University)}\\
Gandhinagar, India}
\and
\IEEEauthorblockN{Abhishek Gupta}
\IEEEauthorblockA{\textit{Mentor} \\
\textit{DAU (Dhirubhai Ambani University)}\\
Gandhinagar, India}
}

\maketitle

%------------------------------------------------------------------
\begin{abstract}
Doctors often make diagnostic mistakes because of cognitive bias, and this is a
major cause of harm to patients. Most virtual patient (VP) simulators only check
whether the student got the right diagnosis. They do not look at how the student
reached it. We built a VP simulator that does look at the process. The student
talks to an AI patient, orders examinations and tests from a fixed list of 27
examinations and 86 investigations that is the same for every case, and then
gives a diagnosis. Our system records every action and uses simple rules to
detect three biases: anchoring, premature closure and confirmation bias. It then
gives feedback about how the student reasoned instead of just telling them the
answer. We tested the system in three ways. First, we checked the bias detectors
on 18 test transcripts that we labelled ourselves. All three detectors caught
every biased transcript (100\% sensitivity), with 94\% accuracy overall. This
testing also found a bug that had made our confirmation-bias detector almost
useless. Second, we describe the size and content of the finished system. Third,
we ran a small pilot with 8 volunteers who did 16 consultations in total.
Between their first and second case, average history coverage went up from
45.3\% to 79.7\% ($p=0.036$) and the average number of questions asked went up
from 6.5 to 9.0 ($p=0.012$). Anchoring dropped from 4 out of 8 students to 0,
and premature closure from 7 to 2, but with only 8 people these changes were not
statistically significant. One student got worse, and one of our measures turned
out to be misleading; we explain both below.
\end{abstract}

\begin{IEEEkeywords}
virtual patient simulator, clinical reasoning, cognitive bias, anchoring,
premature closure, medical education, large language models
\end{IEEEkeywords}

%==================================================================
\section{Introduction}
Making a diagnosis is hard because doctors have to work with incomplete
information. Human thinking has some built-in shortcuts that cause repeated
errors, and these are called cognitive biases. They are a big reason for wrong
diagnoses and avoidable harm~\cite{croskerry2003, graber2005}. Three biases come
up again and again in diagnosis:

\begin{itemize}
  \item \textbf{Anchoring} --- getting stuck on the first idea that comes to mind.
  \item \textbf{Premature closure} --- deciding before gathering enough
  information.
  \item \textbf{Confirmation bias} --- only looking for evidence that supports
  what you already think.
\end{itemize}

Normal assessment only checks the final answer~\cite{norman2010, greengrass2026}.
A student who guesses correctly without asking enough questions still gets marked
correct, even though the same approach could be dangerous with a real patient.
Studies of existing VP systems show that most of them only give this kind of
knowledge feedback and do not check for biased reasoning~\cite{vpsurvey}.

We wanted to close that gap. Our contributions are:

\begin{enumerate}
  \item A working VP simulator that hides which tests are relevant, so students
  cannot guess the diagnosis from the menu.
  \item A simple, rule-based way to detect three biases from what the student
  asked and did.
  \item Testing of those detectors on labelled transcripts, including a bug we
  found and fixed because of that testing.
  \item A small pilot study with 8 volunteers showing how their behaviour
  changed after getting feedback.
\end{enumerate}

%==================================================================
\section{Related Work}
\textbf{Virtual patients using LLMs.}
Recent surveys show that VP systems using large language models are growing
quickly~\cite{vpsurvey}. They are realistic and easy to scale, but they still do
not check the quality of the student's reasoning. In our system the LLM only
plays the patient and writes the feedback text. All the marking is done by our
own Python code.

\textbf{Chatbots in clinical teaching.}
Chatbot teaching tools have been tested before for clinical reasoning skill,
confidence and student satisfaction using before/after
designs~\cite{nursingchatbot}. We used a similar idea and asked students to rate
their confidence before starting each case.

\textbf{Cognitive bias.}
Dual-process theory says we have a fast, automatic way of thinking (System 1)
and a slow, careful one (System 2). The three biases we target are failures of
the fast system~\cite{greengrass2026, norman2010}. Most earlier work is aimed at
working doctors. We aim at students while they are still learning, and use bias
detection as teaching feedback.

\textbf{Explainable feedback.}
Work on explainable AI argues that a system's output should be something a user
can inspect and understand, not a black box~\cite{explainable}. This is why we
used rules instead of a machine learning model. Every flag our system raises can
be traced back to the exact questions the student asked or skipped.

%==================================================================
\section{System Architecture}
The simulator is a web application. The backend is written in Python 3.11 using
Flask, and the frontend is plain HTML, CSS and JavaScript with no framework. The
patient's replies and the feedback text come from an LLM through a
chat-completions API (we used Llama 3.3 70B through Groq; the provider can be
changed in the config file). The code is split into 12 Python files, each with
one job: \texttt{cases.py} holds the clinical content,
\texttt{session\_tracker.py} records what the student does,
\texttt{bias\_detector.py} has the bias rules, \texttt{clinical\_evaluator.py}
marks the diagnosis and tests, \texttt{feedback\_generator.py} writes the
feedback, and \texttt{app.py} handles the web routes.

One design rule we followed strictly: \emph{the LLM never does any marking}. It
only writes text. Every flag, score and verdict is calculated by our own code.
This means the results are repeatable, and we were able to use this later to
check that our collected data was genuine (Section~\ref{sec:pilot}).

\subsection{How a session works}
\begin{enumerate}
  \item \textbf{Before the case.} The student enters a participant ID, their year
  of study, and how confident they feel.
  \item \textbf{History taking.} Every question they type is saved, the question
  counter goes up, and we work out which clinical topics they covered
  (Section~\ref{sec:topic}).
  \item \textbf{Examination and tests.} Everything they click is added to a list
  (no duplicates). The options come from the same universal menu on every case
  (Section~\ref{sec:universal}).
  \item \textbf{Finishing.} When the student submits a diagnosis, the session
  goes to the bias detector and the evaluator, feedback is generated, and the
  whole thing is saved as a JSON file.
\end{enumerate}

\subsection{Saving data and linking sessions}
While a consultation is running, the session is kept in memory on the server. The
browser cookie only holds an ID. When the student finishes, everything is written
to a JSON file with 18 fields. Each file stores the participant ID and a session
number (1 for their first case, 2 for their second), worked out by counting how
many sessions that participant has already done. This is what lets us match a
student's ``before'' and ``after'' sessions. These JSON files, not the web page,
are what we actually analyse, because they keep the whole process and not just
the final answer.

%==================================================================
\section{Clinical Case Design}
\label{sec:cases}
We wrote five cases. Each one has a \emph{trap}: an obvious diagnosis that is
wrong, which a student who jumps to conclusions is likely to pick
(Table~\ref{tab:cases}). Two of the cases cannot be solved from the history
alone, so students have to order the right tests. Every case has 8 required
history topics and a suggested minimum of 7--8 questions.

\begin{table}[t]
\caption{Our five cases and the trap in each one.}
\label{tab:cases}
\centering
\small
\begin{tabular}{@{}p{0.7cm}p{2.4cm}p{2.0cm}p{1.7cm}@{}}
\toprule
\textbf{Case} & \textbf{Presentation} & \textbf{Real diagnosis} & \textbf{Trap} \\
\midrule
1 & Chest pain, 48M & GERD (from NSAIDs) & Cardiac/ACS \\
2 & Breathlessness, 29F & Pulmonary embolism & Chest infection \\
3 & Acute confusion, 74M & UTI/delirium & Stroke/dementia \\
4 & Fatigue, weight gain, 35F & Hypothyroidism & Depression \\
5 & Vomiting, abdominal pain, 21F & DKA (new T1DM) & Gastroenteritis \\
\bottomrule
\end{tabular}
\end{table}

\subsection{The same test list for every case}
\label{sec:universal}
If we only showed the tests that mattered for a case, students could work out the
answer just by reading the menu. So we made one big list that is identical for
all five cases: 27 examinations and 86 investigations. If a test matters for that
case, the student gets the real result. If not, they get a normal result. Behind
the scenes each case marks a few tests as \emph{key}, \emph{reasonable} or
\emph{low value} for marking, but the student never sees these labels.

%==================================================================
\section{How We Detect Bias}
\label{sec:bias}
Each detector gives back a yes/no flag, a score between 0 and 1, a reason written
in plain English, and the evidence it used.

\subsection{Finding topics}
\label{sec:topic}
We check each question against a list of 523 keyword phrases grouped into 40
topics. If a question contains any keyword for a topic, that topic counts as
covered. History coverage is just how many of the case's required topics the
student covered.

\subsection{Anchoring}
Let $q$ be the number of questions, $a$ how many mention the trap diagnosis, and
$m$ how many mention other possibilities. We flag anchoring if either:
\textbf{A1} $q \geq 4$ and $a/q > 0.60$, or \textbf{A2} $a \geq 3$ and $m = 0$.

\subsection{Premature closure}
Let $q_{\min}$ be the minimum questions for that case and $c$ the coverage
fraction. We flag it if either \textbf{P1} $q < q_{\min}$, or \textbf{P2}
$c < 0.60$.

\subsection{Confirmation bias}
Each case has a list of \emph{contradictory clues}: keywords for the evidence
that points away from the trap. We picked these carefully so they do not share
words with the trap keywords. If the student explored $k$ out of $K$ clues, we
flag confirmation bias if either \textbf{C1} they diagnosed the trap and $k = 0$,
or \textbf{C2} $q \geq 5$ and $k/K < 0.25$.

\subsection{Marking the diagnosis}
Separately from the bias detection, we mark the diagnosis as \emph{correct},
\emph{partial}, \emph{anchored} or \emph{other}, and produce a scorecard of which
key examinations and tests were done or skipped. Because these are separate, a
student can be marked correct on the diagnosis and still be flagged for premature
closure at the same time. That combination is exactly what normal marking cannot
show.

%==================================================================
\section{Feedback}
When the student submits their diagnosis, we build a prompt containing what they
did, which biases were flagged, and their scorecard, and ask the LLM for
feedback. The feedback asks reflective questions about what they missed, without
ever using the words ``bias'', ``anchoring'' and so on, because we did not want
to label the student. If the API is down, a backup set of rule-based messages is
used instead, so a student always gets something useful.

%==================================================================
\section{Testing the Detectors}
\label{sec:validation}
\textbf{What we did.} Before trusting the detectors, we wanted to check they
actually work. We wrote 18 test transcripts covering all five cases. Each one is
a realistic set of questions plus a final diagnosis, and we labelled each one
ourselves based on the behaviour it showed, not on what we knew the rules did. We
included four kinds: clearly biased, thorough, quick but open-minded, and a
tricky type where the student is thorough but deliberately avoids the obvious
medical words. Then we ran all 18 through the real system and compared the
output with our labels. The script is \texttt{validate\_detectors.py}.

\textbf{Results.} Across 54 decisions the system was right 94\% of the time
(51 out of 54). All three detectors caught every biased transcript
(Table~\ref{tab:validation}). All three mistakes were false alarms, and all three
happened on the two ``avoiding medical words'' transcripts, which is a direct
result of using keyword matching (see Section~\ref{sec:limitations}).

\begin{table}[t]
\caption{Detector testing on 18 labelled transcripts (54 decisions).}
\label{tab:validation}
\centering
\small
\begin{tabular}{@{}lccc@{}}
\toprule
\textbf{Detector} & \textbf{Sensitivity} & \textbf{Specificity} & \textbf{Accuracy} \\
\midrule
Anchoring          & 100\% & 100\% & 100\% \\
Premature closure  & 100\% & 88\%  & 94\%  \\
Confirmation bias  & 100\% & 82\%  & 89\%  \\
\midrule
\textbf{Overall}   & \multicolumn{3}{c}{94\% (51/54 decisions)} \\
\bottomrule
\end{tabular}
\end{table}

\textbf{A bug this testing found.} Our first version of the confirmation-bias
detector only caught 14\% of the biased transcripts. It missed 6 out of 7. When
we looked into why, we found the problem was that the contradictory clues were
written as full sentences, and those sentences used some of the same words as the
trap keywords. So when a student asked ``any family history of heart problems?''
it matched the clue ``father diabetes no family cardiac heart history'', and the
system thought the student had explored the opposing evidence when they had not.
We rewrote all the clues as short keyword lists with no shared words, and
sensitivity went from 14\% to 100\% without breaking the other two detectors.
This bug was impossible to spot just by reading the code, which is why we think
testing the detectors first was worth doing.

\subsection{Size of the finished system}
Table~\ref{tab:system} shows what we ended up building. The number of keywords
matters because the detectors can only be as good as their word lists.

\begin{table}[t]
\caption{What is in the finished system.}
\label{tab:system}
\centering
\small
\begin{tabular}{@{}lr@{}}
\toprule
\textbf{Item} & \textbf{Count} \\
\midrule
Clinical cases (all with traps) & 5 \\
\quad of which need lab tests & 2 \\
Examinations in the menu & 27 \\
Investigations in the menu & 86 \\
Topics / keyword phrases & 40 / 523 \\
Trap keywords (all cases) & 133 \\
Alternative-diagnosis keywords (all cases) & 134 \\
Contradictory clue sets (all cases) & 30 \\
Fields saved per session & 18 \\
Python files & 12 \\
\bottomrule
\end{tabular}
\end{table}

%==================================================================
\section{Pilot Study}
\label{sec:pilot}
\subsection{What we did}
We asked 8 volunteers (classmates from computer science and medicine, years 2 to
5) to each do two consultations. After the first one they read their feedback,
then they did a second case. Sessions were linked using the participant ID and
session number. We collected 16 consultations in one session on 31 July 2026.
This was an informal test of a teaching prototype. No real patients were
involved and no patient data was used.

\textbf{What we measured}, taken automatically from the saved files by
\texttt{analyze\_sessions.py}: the three bias flags, history coverage, the
fraction of key tests ordered, number of questions asked, and the diagnosis
verdict. For the yes/no flags we used McNemar's exact test and for the numbers we
used the Wilcoxon signed-rank test, since both are made for before/after data
from the same people.

\textbf{Checking the data was real.} Because none of our marking uses the LLM,
we could recompute every bias flag and score from the saved questions and check
it matched what was stored. All 16 sessions matched exactly, all had the full
18-field format, and the question count matched the saved question list in every
file.

\subsection{Results}
Table~\ref{tab:pilot} shows the group results and Table~\ref{tab:perpart} shows
each student separately.

\begin{table}[t]
\caption{Before vs after feedback ($n=8$).}
\label{tab:pilot}
\centering
\small
\begin{tabular}{@{}lccc@{}}
\toprule
\textbf{Measure} & \textbf{Before} & \textbf{After} & \textbf{$p$} \\
\midrule
Anchoring flagged            & 4/8    & 0/8    & 0.125 \\
Premature closure flagged    & 7/8    & 2/8    & 0.125 \\
Confirmation bias flagged    & 5/8    & 2/8    & 0.375 \\
History coverage (average)   & 45.3\% & 79.7\% & \textbf{0.036} \\
Questions asked (average)    & 6.5    & 9.0    & \textbf{0.012} \\
Key tests ordered (average)  & 41.6\% & 20.3\% & 0.123 \\
Correct diagnoses            & 5/8    & 7/8    & --- \\
\bottomrule
\end{tabular}
\end{table}

\begin{table}[t]
\caption{Each student separately. A/P/C = anchoring / premature closure /
confirmation bias. Y means flagged.}
\label{tab:perpart}
\centering
\small
\begin{tabular}{@{}lccccc@{}}
\toprule
\textbf{ID} & \textbf{Yr} & \textbf{Before A/P/C} & \textbf{After A/P/C}
& \textbf{Cov.\ before} & \textbf{Cov.\ after} \\
\midrule
P01 & 3 & Y/Y/Y & --/--/-- & 37.5\% & 100.0\% \\
P02 & 4 & Y/Y/Y & --/--/-- & 37.5\% & 87.5\% \\
P03 & 2 & --/Y/Y & --/--/Y & 25.0\% & 75.0\% \\
P04 & 4 & --/--/-- & --/Y/Y & 75.0\% & 37.5\% \\
P05 & 5 & --/Y/Y & --/--/-- & 75.0\% & 100.0\% \\
P06 & 3 & --/Y/-- & --/--/-- & 50.0\% & 100.0\% \\
P07 & 2 & Y/Y/-- & --/--/-- & 25.0\% & 62.5\% \\
P08 & 3 & Y/Y/Y & --/Y/-- & 37.5\% & 75.0\% \\
\bottomrule
\end{tabular}
\end{table}

Two results were statistically significant. History coverage went up from 45.3\%
to 79.7\% ($p=0.036$), and the number of questions asked went up from 6.5 to 9.0
($p=0.012$). Seven of the eight students covered more history the second time.
All four anchoring flags cleared, and premature closure went from seven students
down to two, but with only 8 people these did not reach significance, so we can
only say the direction looks right.

\textbf{One student got worse.} P04 was the only one who picked up more flags the
second time (premature closure and confirmation bias) and whose coverage went
down, from 75.0\% to 37.5\%, even though they asked more questions
(7~$\rightarrow$~9). P04 also had the best first attempt of anyone. Our guess is
that their second case was simply harder for them, rather than the feedback
making things worse, but since each student only did one pair of cases we cannot
actually tell the difference.

\textbf{One of our measures is misleading.} The key-test number looks like it
dropped a lot, from 41.6\% to 20.3\%, but this is not real. It happens because
different cases need different numbers of key tests. Five of our eight students
went from Case 1 (4 key tests) to Case 2 (8 key tests), so the number they are
being divided by doubles. For example P02 ordered 3 key tests in both sessions:
that is 3/4 = 75\% the first time and 3/8 = 38\% the second time, even though
they did the same amount of work. So this measure cannot be compared across
different cases and we only report it for completeness.

\textbf{Confidence and accuracy.} In our sample, students who felt more confident
before starting did better on their first case: those who rated themselves 4 or 5
got it right 4 times out of 4, and those who rated 1 or 2 got it right 0 times
out of 2. With only 8 students this is just an observation, not a finding.

%==================================================================
\section{Discussion}
Our pilot fits the idea that feedback about the process changes how students
gather information. The two things we measured most directly, how much of the
history they covered and how many questions they asked, both went up
significantly, and all three bias flags moved in the right direction across the
group. The clearest change was history coverage, which makes sense, because the
feedback specifically lists the history areas the student skipped.

There are three things we have to be careful about. First, with only 8 students
none of the yes/no results are significant, so the drop in bias flags is
suggestive but not proven. Second, we had no control group, so we cannot separate
the effect of the feedback from students simply getting better with practice, or
from getting used to the interface. Everyone started with low coverage, so some
improvement on a second try would be expected anyway. Third, one student got
worse, which a group average on its own would have hidden.

The main thing our system does that normal marking cannot is score the diagnosis
and the reasoning separately. Five students got the right diagnosis on their
first case while seven out of eight were flagged for premature closure at the
same time. A normal system would have told most of those students they did well.

\subsection{Limitations}
\label{sec:limitations}
\textbf{Study design.} We had no control group. We also did not keep the case
pair the same for everyone: five students did Case 1 then Case 2, and the other
three did different pairs. We did not swap the order either, so the difficulty of
the second case is mixed up with the effect of the feedback. Our volunteers were
classmates rather than a proper sample of medical students, and some were not
medical students at all. The consultations were also very quick (1.8 minutes on
average), which suggests people went through them faster than they would with a
real patient. We did not give a questionnaire afterwards.

\textbf{Measurement.} Our topic and clue detection uses keyword matching, so a
question worded in an unusual way is not recognised. This caused all three of our
detector mistakes and is the main limit on how accurate our measurements can be.
The thresholds we chose (0.60, 0.60 and 0.25) are sensible but we did not tune
them with data. We labelled the test transcripts ourselves, so having a doctor
label them independently would be stronger. And as explained above, the key-test
measure cannot be compared between cases.

\textbf{System.} Sessions are only saved when the student finishes, so if
someone closes the tab halfway through, that session is lost.

%==================================================================
\section{Conclusion and Future Work}
We built a virtual patient simulator that gives feedback on how a student
reasons instead of just whether they got the answer right. It detects anchoring,
premature closure and confirmation bias using simple rules that can be traced
back to the student's own questions. We tested the detectors on 18 labelled
transcripts and got 100\% sensitivity and 94\% overall accuracy, and that testing
also found and let us fix a bug that had made confirmation-bias detection almost
useless. In a small pilot with 8 students and 16 consultations, history coverage
went up from 45.3\% to 79.7\% ($p=0.036$) and questions asked went up from 6.5 to
9.0 ($p=0.012$) after feedback, with the bias flags moving in the right direction
but not reaching significance with this few people.

What we would do next follows straight from the limitations. To actually claim
that the feedback causes the improvement, we would need a proper study with a
control group who get no feedback, the same pair of cases for everyone with the
order swapped between students, and more participants. On the technical side, the
biggest improvement would be replacing keyword matching with sentence embeddings,
so that questions worded differently are still recognised, since that caused
every mistake our detectors made. We would also like a doctor to label the test
transcripts independently and to tune the thresholds using real data. After that,
we plan to generate new cases automatically from public medical datasets so the
case list can grow without writing each one by hand.

%==================================================================
\section*{Acknowledgment}
We thank Abhishek Gupta for his guidance throughout this internship, the
volunteers who took part in our pilot, and DAU (Dhirubhai Ambani University) for
supporting the project.

%==================================================================
\begin{thebibliography}{00}

\bibitem{croskerry2003}
P.~Croskerry, ``The importance of cognitive errors in diagnosis and strategies
to minimize them,'' \emph{Academic Medicine}, vol.~78, no.~8, pp.~775--780, 2003.

\bibitem{graber2005}
M.~L. Graber, N.~Franklin, and R.~Gordon, ``Diagnostic error in internal
medicine,'' \emph{Archives of Internal Medicine}, vol.~165, no.~13,
pp.~1493--1499, 2005.

\bibitem{norman2010}
G.~R. Norman and K.~W. Eva, ``Diagnostic error and clinical reasoning,''
\emph{Medical Education}, vol.~44, no.~1, pp.~94--100, 2010.

\bibitem{greengrass2026}
C.~J. Greengrass, ``Transforming clinical reasoning---the role of AI in
supporting human cognitive limitations,'' \emph{Frontiers in Digital Health},
vol.~7, art.~1715440, 2026, doi:10.3389/fdgth.2025.1715440.

\bibitem{vpsurvey}
% VERIFY authors/venue/year against the source PDF before submitting.
``A focused survey on patient simulation with large language models,'' 2025.

\bibitem{nursingchatbot}
% VERIFY authors/venue/year against the source PDF before submitting.
``A chatbot education program for nursing students: a randomized study,''
\emph{BMC Medical Education}, 2025.

\bibitem{explainable}
% VERIFY authors/venue/year against the source PDF before submitting.
``Explainable clinical decision reasoning,'' 2024.

\end{thebibliography}

\end{document}
